% Problem record op_9209e4eb15586dec. This TeX file is derived from
% database/problems_json/op_9209e4eb15586dec.json by scripts/sync-tex.mjs; edit the JSON record
% and rerun the script rather than editing this file.
\section{Quantum capacity of a bosonic thermal attenuator}
\label{sec:op_9209e4eb15586dec}

\paragraph{Problem.}
Let $\Phi_{\eta,\nu}$ be the single-mode bosonic thermal attenuator with
transmissivity $0<\eta<1$ and environmental mean photon number
$0<\nu<\infty$, defined by
\begin{equation}
  \Phi_{\eta,\nu}(\rho_A)
  :=\operatorname{Tr}_{E'}\!\left[
     U_\eta(\rho_A\otimes\tau_{\nu,E})U_\eta^\dagger
  \right],
  \qquad
  \tau_{\nu,E}:=\sum_{k=0}^{\infty}
     \frac{\nu^k}{(\nu+1)^{k+1}}
     |k\rangle_E\!\langle k|_E,
  \label{eq:p49-thermal-attenuator}
\end{equation}
where $U_\eta:AE\to BE'$ is a beam-splitter unitary and
$\tau_{\nu,E}$ acts on the environment mode $E$.  For an environment mode of
angular frequency $\omega_E$ at temperature $T$,
$\nu=(e^{\hbar\omega_E/(k_{\rm B}T)}-1)^{-1}$.  Thus $0<T<\infty$ is
equivalent to $0<\nu<\infty$ for fixed $\omega_E>0$; $\nu=0$ corresponds to
$T=0$, while $\nu\to\infty$ as $T\to\infty$.  Its unassisted quantum capacity
is the regularized coherent information
\begin{equation}
  \mathcal Q(\Phi_{\eta,\nu})
  :=\lim_{n\to\infty}\frac1n\sup_{\rho_{A^n}}
       I_{\rm c}(\rho_{A^n},\Phi_{\eta,\nu}^{\otimes n}),
  \qquad
  I_{\rm c}(\rho,\mathcal N)
  :=S(\mathcal N(\rho))-S(\mathcal N^{\rm c}(\rho)),
  \label{eq:p49-quantum-capacity}
\end{equation}
where $S(\sigma):=-\operatorname{Tr}(\sigma\log_2\sigma)$ and
$\mathcal N^{\rm c}$ is any complementary channel.  Determine
Eq.~\eqref{eq:p49-quantum-capacity} for every
$0<\eta<1$ and $0<\nu<\infty$, equivalently every strictly positive finite
environment temperature.

\subsection*{Status}

Unsolved

\subsection*{Source}

Rosati, Mari, and Giovannetti explicitly study the unknown thermal-attenuator
quantum capacity through nonmatching narrow bounds; the exact-capacity
problem is retained after the later non-Gaussian separation
\sourcecite{ref:p49-rosati-mari-giovannetti}{RMG18},
\sourcecite{ref:p49-mele-et-al}{MCF+26}.

\subsection*{Progress}

\begin{itemize}
  \item At zero environmental temperature, the pure-loss capacity is
  \begin{equation}
    \mathcal Q(\Phi_{\eta,0})
    =\left[\log_2\frac{\eta}{1-\eta}\right]_+,
    \qquad [x]_+:=\max\{0,x\}.
    \label{eq:p49-pure-loss-capacity}
  \end{equation}
  Thus Eq.~\eqref{eq:p49-pure-loss-capacity} solves only the boundary case
  $\nu=0$ \sourcecite{ref:p49-wolf-perez-garcia-giedke}{WPG07}.

  \item Optimizing the one-use coherent information over all single-mode
  Gaussian input states gives exactly
  \begin{equation}
    \sup_{\rho\in\mathsf G_1}I_{\rm c}(\rho,\Phi_{\eta,\nu})
    =\left[\log_2\frac{\eta}{1-\eta}-g(\nu)\right]_+,
    \qquad
    g(x):=(x+1)\log_2(x+1)-x\log_2x.
    \label{eq:p49-gaussian-coherent-information}
  \end{equation}
  Holevo and Werner computed the thermal-input limit, Br\'adler proved that
  the thermal-state supremum is attained asymptotically at infinite input
  energy, and Mele et al.\ proved that squeezing cannot improve it within
  $\mathsf G_1$.  Equation~\eqref{eq:p49-gaussian-coherent-information} is
  nevertheless not optimal over arbitrary non-Gaussian inputs
  \sourcecite{ref:p49-holevo-werner}{HW01},
  \sourcecite{ref:p49-bradler}{Br15},
  \sourcecite{ref:p49-mele-et-al}{MCF+26}.

  \item The exact antidegradability region is
  \begin{equation}
    \eta\leq\eta_{\rm AD}(\nu):=
    \frac{\nu+\frac12}{\nu+1},
    \qquad \mathcal Q(\Phi_{\eta,\nu})=0.
    \label{eq:p49-antidegradable-region}
  \end{equation}
  Hence Eq.~\eqref{eq:p49-quantum-capacity} is settled throughout the region
  in Eq.~\eqref{eq:p49-antidegradable-region}
  \sourcecite{ref:p49-lami-et-al}{LKA+19}.

  \item In the non-entanglement-breaking region, a bottleneck decomposition
  gives the analytic upper bound
  \begin{equation}
    \mathcal Q(\Phi_{\eta,\nu})
    \leq\left[
      \log_2\frac{\eta-(1-\eta)\nu}{(1-\eta)(\nu+1)}
    \right]_+
    \quad\text{when }\eta>(1-\eta)\nu.
    \label{eq:p49-bottleneck-bound}
  \end{equation}
  The channel is entanglement breaking, and therefore has zero capacity,
  when the displayed condition fails.  The bound in
  Eq.~\eqref{eq:p49-bottleneck-bound} was obtained independently by Rosati
  et al.\ and Noh et al.\
  \sourcecite{ref:p49-rosati-mari-giovannetti}{RMG18},
  \sourcecite{ref:p49-noh-albert-jiang}{NAJ19}.  Degradable extensions and
  channel-concatenation orderings subsequently produced tighter
  parameter-dependent upper bounds, but no known upper bound is tight
  throughout the non-antidegradable region
  \sourcecite{ref:p49-fanizza-kianvash-giovannetti}{FKG21},
  \sourcecite{ref:p49-kianvash-fanizza-giovannetti}{KFG24}.

  \item Asymptotic multimode symplectic-lattice GKP codes achieve a rate whose
  gap from the upper bound in Eq.~\eqref{eq:p49-bottleneck-bound} is at most
  $\log_2 e\simeq1.443$ qubits per channel use in the energy-unconstrained
  setting.  This is a constant-gap coding theorem, not an exact capacity
  formula \sourcecite{ref:p49-noh-albert-jiang}{NAJ19}.

  \item The latest rigorous lower bounds use non-Gaussian one-mode inputs:
  \begin{equation}
    \begin{aligned}
      \sup_{\rho\in\mathsf G_1}I_{\rm c}(\rho,\Phi_{4/5,1})
        &=0,
      &\qquad
      \mathcal Q(\Phi_{4/5,1})
        &\geq4.7\times10^{-4},\\
      \mathcal Q(\Phi_{0.7841,1})
        &\geq1.4312\times10^{-7},\\
      \eta_{\rm c}
        &:=\inf\{\eta:\mathcal Q(\Phi_{\eta,1})>0\},
      &\qquad
      0.75\leq\eta_{\rm c}
        &\leq0.7841.
    \end{aligned}
    \label{eq:p49-nongaussian-separation}
  \end{equation}
  The first certificate uses an explicit rank-two state supported on six
  Fock levels; numerical optimization over related families reaches about
  $8.4\times10^{-3}$ qubits per use at $(\eta,\nu)=(0.8,1)$, but that larger
  value is not certified.  The second certificate in
  Eq.~\eqref{eq:p49-nongaussian-separation} also propagates by data processing
  to a two-dimensional parameter region.  These results disprove
  single-mode Gaussian optimality but do not determine the capacity
  \sourcecite{ref:p49-mele-et-al}{MCF+26}.
\end{itemize}

\subsection*{References}

\begin{enumerate}[label={},leftmargin=4.8em,labelsep=0.5em]
  \footnotesize
  \item[\textup{[WPG07]}]\label{ref:p49-wolf-perez-garcia-giedke}
  M. M. Wolf, D. P\'erez-Garc\'ia, and G. Giedke,
  ``Quantum Capacities of Bosonic Channels,''
  \emph{Physical Review Letters} \textbf{98}, 130501 (2007).
  \href{https://doi.org/10.1103/PhysRevLett.98.130501}{doi:10.1103/PhysRevLett.98.130501};
  \href{https://arxiv.org/abs/quant-ph/0606132}{arXiv:quant-ph/0606132}.

  \item[\textup{[HW01]}]\label{ref:p49-holevo-werner}
  A. S. Holevo and R. F. Werner,
  ``Evaluating Capacities of Bosonic Gaussian Channels,''
  \emph{Physical Review A} \textbf{63}, 032312 (2001).
  \href{https://doi.org/10.1103/PhysRevA.63.032312}{doi:10.1103/PhysRevA.63.032312};
  \href{https://arxiv.org/abs/quant-ph/9912067}{arXiv:quant-ph/9912067}.

  \item[\textup{[Br15]}]\label{ref:p49-bradler}
  K. Br\'adler,
  ``Coherent Information of One-Mode Gaussian Channels---The General Case
  of Non-Zero Added Classical Noise,''
  \emph{Journal of Physics A: Mathematical and Theoretical} \textbf{48},
  125301 (2015).
  \href{https://doi.org/10.1088/1751-8113/48/12/125301}{doi:10.1088/1751-8113/48/12/125301};
  \href{https://arxiv.org/abs/1406.6110}{arXiv:1406.6110}.

  \item[\textup{[RMG18]}]\label{ref:p49-rosati-mari-giovannetti}
  M. Rosati, A. Mari, and V. Giovannetti,
  ``Narrow Bounds for the Quantum Capacity of Thermal Attenuators,''
  \emph{Nature Communications} \textbf{9}, 4339 (2018).
  \href{https://doi.org/10.1038/s41467-018-06848-0}{doi:10.1038/s41467-018-06848-0};
  \href{https://arxiv.org/abs/1801.04731}{arXiv:1801.04731}.

  \item[\textup{[NAJ19]}]\label{ref:p49-noh-albert-jiang}
  {\sloppy
  K. Noh, V. V. Albert, and L. Jiang,
  ``Quantum Capacity Bounds of Gaussian Thermal Loss Channels and
  Achievable Rates with Gottesman--Kitaev--Preskill Codes,''
  \emph{IEEE Transactions on Information Theory} \textbf{65}, 2563--2582
  (2019).
  \href{https://doi.org/10.1109/TIT.2018.2873764}{doi:10.1109/TIT.2018.2873764};
  \href{https://arxiv.org/abs/1801.07271}{arXiv:1801.07271}.
  \par}

  \item[\textup{[LKA+19]}]\label{ref:p49-lami-et-al}
  L. Lami, S. Khatri, G. Adesso, and M. M. Wilde,
  ``Extendibility of Bosonic Gaussian States,''
  \emph{Physical Review Letters} \textbf{123}, 050501 (2019).
  \href{https://doi.org/10.1103/PhysRevLett.123.050501}{doi:10.1103/PhysRevLett.123.050501};
  \href{https://arxiv.org/abs/1904.02692}{arXiv:1904.02692}.

  \item[\textup{[FKG21]}]\label{ref:p49-fanizza-kianvash-giovannetti}
  M. Fanizza, F. Kianvash, and V. Giovannetti,
  ``Estimating Quantum and Private Capacities of Gaussian Channels via
  Degradable Extensions,''
  \emph{Physical Review Letters} \textbf{127}, 210501 (2021).
  \href{https://doi.org/10.1103/PhysRevLett.127.210501}{doi:10.1103/PhysRevLett.127.210501};
  \href{https://arxiv.org/abs/2103.09569}{arXiv:2103.09569}.

  \item[\textup{[KFG24]}]\label{ref:p49-kianvash-fanizza-giovannetti}
  F. Kianvash, M. Fanizza, and V. Giovannetti,
  ``Low-Ground/High-Ground Capacity Regions Analysis for Bosonic Gaussian
  Channels,''
  \emph{International Journal of Quantum Information} \textbf{22}, 2440005
  (2024).
  \href{https://doi.org/10.1142/S0219749924400057}{doi:10.1142/S0219749924400057};
  \href{https://arxiv.org/abs/2306.16350}{arXiv:2306.16350}.

  \item[\textup{[MCF+26]}]\label{ref:p49-mele-et-al}
  F. A. Mele, G. Catalano, M. Fanizza, V. Giovannetti, and L. Lami,
  ``Bosonic Quantum Communication Beyond the Thermal Threshold,''
  arXiv preprint (2026).
  \href{https://arxiv.org/abs/2607.27449}{arXiv:2607.27449}.
\end{enumerate}

\subsection*{Comment}

For generic $\nu>0$ outside the zero-capacity regions, the lower and upper
bounds do not coincide.  Non-Gaussian one-mode optimization and possible
multimode superadditivity both remain relevant to
Eq.~\eqref{eq:p49-quantum-capacity}.

\subsection*{Field}

Quantum Shannon theory

\subsection*{Topic}

Bosonic channels; Quantum capacity; Coherent information; Gaussian quantum information; Continuous-variable systems

\subsection*{ID}

\texttt{op\_9209e4eb15586dec}
